% The Figure structure intentionally includes the non-floating caption.
\tagstructbegin{tag=Figure,alttext={Key 与 Value 的存储、查询和按权重读取}}
% Caption first in the PDF/UA-2 Figure structure; remaining prose follows.
\begin{center}\captionof{figure}{Key 的按行存放与转置：两个视图是一组相同数据}\label{fig:key-layout}\end{center}
\begin{tcolorbox}[
  enhanced,breakable=false,colback=paper,colframe=linegray,
  boxrule=.7pt,arc=2mm,left=3mm,right=3mm,top=3mm,bottom=3mm,
]
\small
\noindent{\sffamily\bfseries\color{ink}图解：为什么一会儿写 $K$，一会儿写 $K^\top$？}\par\smallskip
沿用本章开头的 $K$：$T_k=3$ 条 Key，每条 $d_k=2$ 个数字。
\textbf{相同底色表示同一条 Key：两栏数字相同，只是摆放方向不同。}

\vspace{1mm}
\noindent
\begin{minipage}[t]{.49\linewidth}
\centering
\textbf{列向量 Query 的写法}\\[1mm]
\begin{tikzpicture}[x=10mm,y=7mm]
\foreach \r/\shade in {1/paper,2/linegray!25,3/linegray!15}{
  \foreach \c in {1,2}{
    \filldraw[draw=linegray,fill=\shade,line width=.55pt]
      (\c-1,1-\r) rectangle (\c,-\r);
  }
}
\node at (.5,-.5){1}; \node at (1.5,-.5){0};
\node at (.5,-1.5){1}; \node at (1.5,-1.5){1};
\node at (.5,-2.5){2}; \node at (1.5,-2.5){1};
\end{tikzpicture}
\\[.5mm]
$K\in\R^{T_k\times d_k}$，{\footnotesize 每行一条 Key}\\[.5mm]
$K\vect q=\vect s$
\end{minipage}%
\begin{minipage}[t]{.49\linewidth}
\centering
\textbf{行向量 Query 的写法}\\[1mm]
\begin{tikzpicture}[x=10mm,y=7mm]
\foreach \c/\shade in {1/paper,2/linegray!25,3/linegray!15}{
  \foreach \r in {1,2}{
    \filldraw[draw=linegray,fill=\shade,line width=.55pt]
      (\c-1,1-\r) rectangle (\c,-\r);
  }
}
\node at (.5,-.5){1}; \node at (1.5,-.5){1}; \node at (2.5,-.5){2};
\node at (.5,-1.5){0}; \node at (1.5,-1.5){1}; \node at (2.5,-1.5){1};
\end{tikzpicture}
\\[.5mm]
$K^\top\in\R^{d_k\times T_k}$，{\footnotesize 每列一条 Key}\\[.5mm]
$\vect q^\top K^\top=\vect s^\top$
\end{minipage}

\par\smallskip
\noindent\textbf{形状：}$\vect q^\top:1\times d_k$ 乘 $K^\top:d_k\times T_k$ 得 $1\times T_k$ 的分数；
两种写法互为\textbf{整体转置}：{\color{flowcolor}\bfseries $(K\vect q)^\top=\vect q^\top K^\top$}。

\noindent\textbf{Value 同理：}$V:T_k\times d_v$ 每行一条 Value，
$\vect a^\top V=\vect y^\top$（$1\times d_v$）与 $V^\top\vect a=\vect y$ 是同一次读取；
$d_v$ 不必等于 $d_k$。

\noindent\textbf{多个 Query：}写成 $QK^\top$；它与 $KQ^\top$ 互为转置，不能随意交换。
\par % Flush the final box paragraph before the Figure caption tag.
\end{tcolorbox}

\tagstructend
